\clearpage\pagestyle{otMain}\setlength{\otBlocEspacement}{2.5mm}\fontsize{11}{15}\selectfont\corrigesChapitre{3}{Puissances et notation scientifique}
\begin{blocLivre}{cor-3-1}\begin{corrige}{1}{Déplier une puissance}\demarche{Compter les facteurs.}
$5^3=5\times5\times5=\res{125}$. Il y a trois facteurs 5 ; $5\times3$ donnerait 15 et ne représente pas cette puissance.
\end{corrige}\end{blocLivre}
\begin{blocLivre}{cor-3-2}\begin{corrige}{2}{Attention aux parenthèses}\demarche{Repérer la base exacte de la puissance.}
$(-3)^2=(-3)\times(-3)=\res9$. $-3^2=-(3\times3)=\res{-9}$. Enfin $(-3)^3=\res{-27}$, car trois facteurs négatifs donnent un produit négatif.
\end{corrige}\end{blocLivre}
\begin{blocLivre}{cor-3-3}\begin{corrige}{3}{Puissances de dix}\demarche{Relier exposant négatif et inverse.}
$10^5=\res{100\,000}$ ; $10^{-4}=1/10\,000=\res{0{,}0001}$ ; $10^0=\res1$. L’exposant négatif donne l’inverse d’une puissance positive de 10.
\end{corrige}\end{blocLivre}
\begin{blocLivre}{cor-3-4}\begin{corrige}{4}{Vers la notation scientifique}\demarche{Placer la virgule après le premier chiffre non nul.}
$7\,240\,000=\res{7{,}24\times10^6}$ et $0{,}000091=\res{9{,}1\times10^{-5}}$. Les coefficients 7,24 et 9,1 sont compris entre 1 inclus et 10 exclu.
\end{corrige}\end{blocLivre}
\begin{blocLivre}{cor-3-5}\begin{corrige}{5}{Lire un préfixe}\demarche{Remplacer chaque préfixe par sa puissance de 10.}
Le préfixe kilo signifie $10^3$ : $6$ km $=\res{6000\ \mathrm m}$. Micro signifie $10^{-6}$ : $8$ microlitres $=\res{8\times10^{-6}\ \mathrm L}$, soit 0,000008 L.
\end{corrige}\end{blocLivre}
\begin{blocLivre}{cor-3-6}\begin{corrige}{6}{Comparer des petites valeurs}\demarche{Passer si nécessaire au même exposant.}
Au même exposant : $4{,}8\times10^{-5}=48\times10^{-6}$. Comme $48>7{,}2$, le premier nombre est le plus grand. Un exposant $-5$ est supérieur à $-6$.
\end{corrige}\end{blocLivre}
\begin{blocLivre}{cor-3-7}\begin{corrige}{7}{Retrouver les facteurs dix}\demarche{Revenir aux produits et simplifier.}
Le produit contient sept facteurs 10 : $\res{10^7}$. Dans le quotient, deux facteurs se simplifient ; il reste $1/10^3=\res{10^{-3}}$.
\end{corrige}\end{blocLivre}
\begin{blocLivre}{cor-3-8}\begin{corrige}{8}{Normaliser une écriture}\demarche{Compenser tout déplacement de la virgule.}
$48=4{,}8\times10$, donc $\res{4{,}8\times10^7}$. $0{,}62=6{,}2/10$, donc $\res{6{,}2\times10^{-4}}$. Le nombre conserve la même valeur.
\end{corrige}\end{blocLivre}
\begin{blocLivre}{cor-3-9}\begin{corrige}{9}{Multiplier deux écritures}\demarche{Regrouper coefficients et puissances.}
Les coefficients donnent 6. Les facteurs de 10 donnent $10^4/10^6=10^{-2}$. Résultat : $\res{6\times10^{-2}=0{,}06}$.
\end{corrige}\end{blocLivre}
\begin{blocLivre}{cor-3-10}\begin{corrige}{10}{Diviser et estimer}\demarche{Diviser les coefficients et simplifier les facteurs 10.}
Le quotient des coefficients vaut $8{,}4/2=4{,}2$. Après simplification de trois facteurs 10, il reste $10^4$. Résultat : $\res{4{,}2\times10^4=42\,000}$.
\end{corrige}\end{blocLivre}
\begin{blocLivre}{cor-3-11}\begin{corrige}{11}{Des fichiers}\demarche{Choisir la même unité avant de multiplier.}
$2400\times2{,}5=\res{6000}$ Mo. Comme 1 Go = 1000 Mo, le total vaut $\res6$ Go. On peut aussi obtenir $6\times10^9$ octets.
\end{corrige}\end{blocLivre}
\begin{blocLivre}{cor-3-12}\begin{corrige}{12}{Des longueurs microscopiques}\demarche{Convertir les deux mesures dans une seule unité.}
$0{,}8$ micromètre $=800$ nm car 1 micromètre vaut 1000 nm. Le second élément est plus long de $800-450=\res{350}$ nm. En mètres, les tailles sont $4{,}5\times10^{-7}$ et $8\times10^{-7}$.
\end{corrige}\end{blocLivre}
\begin{blocLivre}{cor-3-13}\begin{corrige}{13}{Un ordre de grandeur}\demarche{Estimer avant de calculer exactement.}
En arrondissant $3{,}8$ à 4, on prévoit environ $4\times10^5/(2\times10^3)=200$ étapes. Le calcul exact donne $380\,000/2000=\res{190}$. L’estimation est cohérente.
\end{corrige}\end{blocLivre}
\begin{blocLivre}{cor-3-14}\begin{corrige}{14}{Une répétition qui double}\demarche{Distinguer valeur initiale et facteur répété.}
Chaque étape ajoute un facteur 2. Après six étapes : $3\times2^6=3\times64=\res{192}$. Le facteur initial 3 n’est pas la base de la puissance.
\end{corrige}\end{blocLivre}
\begin{blocLivre}{cor-3-15}\begin{corrige}{15}{Une somme à ne pas confondre}\demarche{Identifier l’opération principale.}
Pour la somme, $8\times10^4=0{,}8\times10^5$, donc $A=\res{4\times10^5}$. Pour le produit, $3{,}2\times8=25{,}6$ et neuf facteurs 10 donnent $B=25{,}6\times10^9=\res{2{,}56\times10^{10}}$. Une somme ne se traite pas comme un produit.
\end{corrige}\end{blocLivre}
